% GENERATED FILE. Edit canonical lessons, then run npm run generate:books.
% canonical-source-hash: fnv1a64:45adf6d4321c46a2
% canonical-lessons: MW-C264-gharno, MW-C264-luhar, MW-C264-thokno, MW-C264-kumhar, MW-C264-suthar

\chapter{Smith, Potter and Carpenter}
\label{ch:mw-smith-potter-and-carpenter}
\hlchaptermodality{264}

\begin{chapteropening}
\textbf{By the end of this chapter:} \emph{I can say what the smith, the potter and the carpenter make.}

Complete the last lesson of chapter 264: I can say what the smith, the potter and the carpenter make.
\end{chapteropening}

\section[\texorpdfstring{\mw{घड़णो} (ghaṛṇo)}{घड़णो (ghaṛṇo)}]{\mw{घड़णो} (ghaṛṇo) — to shape; to forge (metal or clay)}

\label{lesson:MW-C264-gharno}



\begin{quote}
Before the new one: say the Marwadi for a washerman, then the Marwadi for a goldsmith.
\end{quote}

\subsection*{You'll want to know: \mw{घड़णो}}
\textbf{\mw{घड़णो}} — \emph{ghaṛṇo} — \textquotedblleft{}to shape; to forge (metal or clay)\textquotedblright{}.

\begin{quote}
\mw{सुनार} \mw{अंगूठी} \mw{घड़ै} \mw{है।}
\end{quote}

\begin{quote}
\emph{sunār aṅgūṭhī ghaṛai hai.}
\end{quote}

\begin{quote}
The goldsmith shapes a ring.
\end{quote}

Past (I …): \mw{म्हैं} \mw{कड़ो} \mw{घड़्यो}, \emph{mhaĩ kaṛo ghaṛyo}.

\begin{cousinweb}
\textbf{\mw{घड़णो}} (\emph{ghaṛṇo}) is inherited from Sanskrit \emph{ghaṭate}, \textquotedblleft{}is busy with\textquotedblright{}; CDIAL lists Old Marwari \emph{ghaṛī}, \textquotedblleft{}formed\textquotedblright{}, under it (CDIAL 4407).
\end{cousinweb}

\begin{grammarlens}[title={What you've built}]
One more word for this chapter.
\end{grammarlens}

\subsection*{Your turn}
\begin{itemize}
\raggedright
  \item \emph{Say it:} \emph{ghaṛṇo}
  \item \emph{Say it:} \emph{ghaṛṇo}, once more
  \item \emph{Say it:} \emph{sunār aṅgūṭhī ghaṛai hai.} — The goldsmith shapes a ring.
  \item \emph{Recall:} say the Marwadi for a washerman, then the Marwadi for a goldsmith, then say \emph{ghaṛṇo} again
\end{itemize}

\begin{tcolorbox}[breakable,colback=teal!4,colframe=teal!35!black,title={Before you move on}]
What does \mw{घड़णो} mean? (\textquotedblleft{}To shape; to forge (metal or clay)\textquotedblright{}.) Say it once more.
\end{tcolorbox}
\section[\texorpdfstring{\mw{लुहार} (luhār)}{लुहार (luhār)}]{\mw{लुहार} (luhār) — a blacksmith}

\label{lesson:MW-C264-luhar}



\begin{quote}
Before the new one: say the Marwadi for a goldsmith, then the Marwadi for to shape (metal or clay).
\end{quote}

\subsection*{You'll want to know: \mw{लुहार}}
\textbf{\mw{लुहार}} — \emph{luhār} — \textquotedblleft{}a blacksmith\textquotedblright{}.

\begin{quote}
\mw{लुहार} \mw{कुल्हाड़ी} \mw{बणावै} \mw{है।}
\end{quote}

\begin{quote}
\emph{luhār kulhāṛī baṇāvai hai.}
\end{quote}

\begin{quote}
The blacksmith makes an axe.
\end{quote}

A masculine noun.

\begin{cousinweb}
\textbf{\mw{लुहार}} (\emph{luhār}) is inherited from Sanskrit \emph{lohakāra}, \textquotedblleft{}blacksmith\textquotedblright{}. CDIAL gives no Marwari form in this entry; it lists Hindi \emph{luhār} (CDIAL 11159).
\end{cousinweb}

\begin{grammarlens}[title={What you've built}]
One more word for this chapter.
\end{grammarlens}

\subsection*{Your turn}
\begin{itemize}
\raggedright
  \item \emph{Say it:} \emph{luhār}
  \item \emph{Say it:} \emph{luhār}, once more
  \item \emph{Say it:} \emph{luhār kulhāṛī baṇāvai hai.} — The blacksmith makes an axe.
  \item \emph{Recall:} say the Marwadi for a goldsmith, then the Marwadi for to shape (metal or clay), then say \emph{luhār} again
\end{itemize}

\begin{tcolorbox}[breakable,colback=teal!4,colframe=teal!35!black,title={Before you move on}]
What does \mw{लुहार} mean? (\textquotedblleft{}A blacksmith\textquotedblright{}.) Say it once more.
\end{tcolorbox}
\section[\texorpdfstring{\mw{ठोकणो} (ṭhokṇo)}{ठोकणो (ṭhokṇo)}]{\mw{ठोकणो} (ṭhokṇo) — to hammer; to knock in}

\label{lesson:MW-C264-thokno}



\begin{quote}
Before the new one: say the Marwadi for to shape (metal or clay), then the Marwadi for a blacksmith.
\end{quote}

\subsection*{You'll want to know: \mw{ठोकणो}}
\textbf{\mw{ठोकणो}} — \emph{ṭhokṇo} — \textquotedblleft{}to hammer; to knock in\textquotedblright{}.

\begin{quote}
\mw{लुहार} \mw{खीलो} \mw{ठोकै} \mw{है।}
\end{quote}

\begin{quote}
\emph{luhār khīlo ṭhokai hai.}
\end{quote}

\begin{quote}
The blacksmith hammers in a nail.
\end{quote}

Past (I …): \mw{म्हैं} \mw{खीलो} \mw{ठोक्यो}, \emph{mhaĩ khīlo ṭhokyo}.

\begin{cousinweb}
\textbf{\mw{ठोकणो}} (\emph{ṭhokṇo}) is inherited from a reconstructed \emph{*ṭhōkk}, \textquotedblleft{}knock\textquotedblright{}. CDIAL gives no Marwari form in this entry; it lists Hindi \emph{ṭhoknā}, \textquotedblleft{}to knock, make firm\textquotedblright{} (CDIAL 5513).
\end{cousinweb}

\begin{grammarlens}[title={What you've built}]
One more word for this chapter.
\end{grammarlens}

\subsection*{Your turn}
\begin{itemize}
\raggedright
  \item \emph{Say it:} \emph{ṭhokṇo}
  \item \emph{Say it:} \emph{ṭhokṇo}, once more
  \item \emph{Say it:} \emph{luhār khīlo ṭhokai hai.} — The blacksmith hammers in a nail.
  \item \emph{Recall:} say the Marwadi for to shape (metal or clay), then the Marwadi for a blacksmith, then say \emph{ṭhokṇo} again
\end{itemize}

\begin{tcolorbox}[breakable,colback=teal!4,colframe=teal!35!black,title={Before you move on}]
What does \mw{ठोकणो} mean? (\textquotedblleft{}To hammer; to knock in\textquotedblright{}.) Say it once more.
\end{tcolorbox}
\section[\texorpdfstring{\mw{कुम्हार} (kumhār)}{कुम्हार (kumhār)}]{\mw{कुम्हार} (kumhār) — a potter}

\label{lesson:MW-C264-kumhar}



\begin{quote}
Before the new one: say the Marwadi for a blacksmith, then the Marwadi for to hammer.
\end{quote}

\subsection*{You'll want to know: \mw{कुम्हार}}
\textbf{\mw{कुम्हार}} — \emph{kumhār} — \textquotedblleft{}a potter\textquotedblright{}.

\begin{quote}
\mw{कुम्हार} \mw{घड़ो} \mw{बणावै} \mw{है।}
\end{quote}

\begin{quote}
\emph{kumhār ghaṛo baṇāvai hai.}
\end{quote}

\begin{quote}
The potter makes a water pot.
\end{quote}

A masculine noun.

\begin{cousinweb}
\textbf{\mw{कुम्हार}} (\emph{kumhār}) is inherited from Sanskrit \emph{kumbhakāra}, \textquotedblleft{}potter\textquotedblright{}. CDIAL gives no Marwari form in this entry; it lists Hindi \emph{kumhār} (CDIAL 3310).
\end{cousinweb}

\begin{grammarlens}[title={What you've built}]
One more word for this chapter.
\end{grammarlens}

\subsection*{Your turn}
\begin{itemize}
\raggedright
  \item \emph{Say it:} \emph{kumhār}
  \item \emph{Say it:} \emph{kumhār}, once more
  \item \emph{Say it:} \emph{kumhār ghaṛo baṇāvai hai.} — The potter makes a water pot.
  \item \emph{Recall:} say the Marwadi for a blacksmith, then the Marwadi for to hammer, then say \emph{kumhār} again
\end{itemize}

\begin{tcolorbox}[breakable,colback=teal!4,colframe=teal!35!black,title={Before you move on}]
What does \mw{कुम्हार} mean? (\textquotedblleft{}A potter\textquotedblright{}.) Say it once more.
\end{tcolorbox}
\section[\texorpdfstring{\mw{सुथार} (suthār)}{सुथार (suthār)}]{\mw{सुथार} (suthār) — a carpenter}

\label{lesson:MW-C264-suthar}



\begin{quote}
Before the new one: say the Marwadi for to hammer, then the Marwadi for a potter.
\end{quote}

\subsection*{You'll want to know: \mw{सुथार}}
\textbf{\mw{सुथार}} — \emph{suthār} — \textquotedblleft{}a carpenter\textquotedblright{}.

\begin{quote}
\mw{सुथार} \mw{किवाड़} \mw{बणावै} \mw{है।}
\end{quote}

\begin{quote}
\emph{suthār kivāṛ baṇāvai hai.}
\end{quote}

\begin{quote}
The carpenter makes a door.
\end{quote}

A masculine noun.

\begin{cousinweb}
\textbf{\mw{सुथार}} (\emph{suthār}) is inherited from Sanskrit \emph{sūtradhāra}, \textquotedblleft{}the holder of the measuring line\textquotedblright{}, a carpenter. CDIAL gives no Marwari form in this entry; it lists Gujarati \emph{suthār} and Hindi \emph{sutār} (CDIAL 13563).
\end{cousinweb}

\begin{grammarlens}[title={What you've built}]
One more word for this chapter.
\end{grammarlens}

\subsection*{Your turn}
\begin{itemize}
\raggedright
  \item \emph{Say it:} \emph{suthār}
  \item \emph{Say it:} \emph{suthār}, once more
  \item \emph{Say it:} \emph{suthār kivāṛ baṇāvai hai.} — The carpenter makes a door.
  \item \emph{Recall:} say the Marwadi for to hammer, then the Marwadi for a potter, then say \emph{suthār} again
\end{itemize}

\begin{tcolorbox}[breakable,colback=teal!4,colframe=teal!35!black,title={Before you move on}]
What does \mw{सुथार} mean? (\textquotedblleft{}A carpenter\textquotedblright{}.) Say it once more.
\end{tcolorbox}
